\section{Solving the free-series equations}
\label{sec:free-correction}

We now use the response and cyclic-sum results to keep the output of a
substitution conjugate to its constant term.  The input retains two prescribed
linear coefficients.  This is the point at which free words, rather than
commuting parameters, enter the construction essentially.

Work over a common split scalar field $E$ in which the preceding evaluations
and spectral projections are defined.  All series substitutions below have
small constant terms and are interpreted by the coefficientwise valuation
convergence established earlier.  In particular, the letters of a free series
commute with its matrix coefficients; the letters do not commute with each
other.

Put $m=11$ and index the blocks by $i\in\mathbb Z/m\mathbb Z$.
Let $P_i$ be jointly independent small generic pairs of $b\times b$
matrices.  Set $R=\operatorname{diag}(P_i)$, a pair of $mb\times mb$
matrices.  The old coefficient algebra acts on a separate factor of dimension
$d$.  Write $\lambda_{i,h}$ and $\pi_{i,h}$ for the distinct eigenvalues
and spectral projections of $w(P_i)$.  Eigenvalues belonging to different
blocks are distinct.  For a substitution $Z$ with constant term $R$, let
$S_{i,h}(Z)$ denote the fixed-root Schur expression at $\lambda_{i,h}$,
using the constant projection $\pi_{i,h}$ in the $i$-th block.

We recall precisely the response information needed here.  In each block
choose one eigenvalue label from every part of the active partition, giving
$a$ selected labels.  The selected loop-response map is
\[
 D_i:M_b(E)^2\longrightarrow E^a.
\]
Put $L_i=\operatorname{im}D_i$ and $H_i=E^a/L_i$.
Proposition~\ref{prop:compatible-relations} supplies a vector space
$G=\mathcal G$ of dimension $g$ and coordinate isomorphisms
\[
 p_i:G\longrightarrow H_i^*.
\]
For a generic fixed input arrow $U$ on $i\leftarrow j$, varying its
opposite arrow $V$ gives two selected responses
\[
 m_i(U,V)\in E^a,
 \qquad m_j(V,U)\in E^a.
\]
The annihilator of these paired responses in $H_i^*\oplus H_j^*$ is
exactly
\begin{equation}
 \{(p_i(c),p_j(c)):c\in G\}.
 \label{eq:correction-paired-annihilator}
\end{equation}
The identifications $p_i$ are compatible for the whole collection of blocks
and for all the generic arrows used below.

The transfer result supplies the following additional facts.  If all Schur
expressions vanish in degrees less than $\ell$, their degree-$\ell$
coefficients are scalar on the projected eigenspaces.  The resulting tuples
belong to their active spaces and are determined by their selected entries.
Denote the selected error for a word $\omega$ of length $\ell$ in block
$i$ by $e_{i,\omega}\in E^a$.  The cyclic-sum result says that, for every
orbit $\mathcal O$ of words under cyclic rotation,
\begin{equation}
 \sum_{\omega\in\mathcal O}\sum_i
 p_i(c)\bigl(e_{i,\omega}+L_i\bigr)=0
 \qquad(c\in G).
 \label{eq:correction-cyclic-input}
\end{equation}
These statements hold for every partial substitution under consideration,
including after an extension of its scalar field.  The preceding results
therefore supply all the hypotheses used in the next construction.

\begin{proposition}
\label{prop:eleven-block-correction}
Prescribe generic input-arrow pairs on the edges
\[
 \begin{array}{c|c}
 \text{letter}&\text{offsets }i-j\text{ of edges }i\leftarrow j\\ \hline
 s_1&1,2\\
 s_2&3,4.
 \end{array}
\]
Under the response and cyclic-sum hypotheses above, there is a pair
$Z(s_1,s_2)\in M_{mb}(E)\langle\!\langle s_1,s_2\rangle\!\rangle^2$
with constant term $R$ and exactly these prescribed linear coefficients,
for which
\[
 S_{i,h}(Z)=0\qquad\text{for every }i,h.
\]
It can be constructed by finite linear systems.  At step $\ell\ge3$ the
new unknowns are coefficients of degree $\ell-1$ on the reverse edges
and coefficients of degree $\ell$ on the diagonal blocks.
\end{proposition}

\begin{proof}
None of the four forward offsets is zero.  The sum of two forward offsets
lies between $2$ and $8$, and is therefore nonzero modulo $11$.
Thus a path of one or two forward steps cannot return to its initial block.
Constant input steps preserve each block.  The path expansions for the
output and for its Schur expressions consequently show that the prescribed
constant and linear terms, with zero higher coefficients, solve all the
Schur equations through degree two.

Suppose a partial substitution solves them in degrees less than
$\ell\ge3$.  Add a correction $V B$ on a reverse edge $j\leftarrow i$,
where $B$ is a word of length $\ell-1$, and add arbitrary degree-$\ell$
loop corrections.  At the block-diagonal constant term, an off-diagonal
input variation gives an output variation in the same off-diagonal block.
Its compression on any fixed root block is zero.  Thus the reverse
correction changes no Schur equation through degree $\ell-1$.

Its contribution in degree $\ell$ must pair it with a single prescribed
linear arrow.  There is a unique forward edge that returns along this
reverse edge: the reverse offsets $10,9,8,7$ correspond respectively to
$1,2,3,4$.  If that edge has coefficient $U$ and color $s_c$, the
contributions are
\begin{equation}
 \begin{array}{c|c|c}
 \text{block}&\text{word}&\text{selected response}\\ \hline
 i&s_cB&m_i(U,V)\\
 j&Bs_c&m_j(V,U).
 \end{array}
 \label{eq:correction-leading-pair}
\end{equation}
This is exactly the two-arrow Schur calculation defining the paired
responses.  A previously present coefficient of degree at least two
cannot interact with the new degree-$\ell-1$ coefficient in degree
$\ell$.  Two new corrections have total degree at least
$2(\ell-1)>\ell$.  The degree-$\ell$ variation is therefore linear.
Loop corrections give $D_i$ independently at every vertex and every
word of length $\ell$.

We compute the annihilator of this finite correction map.  Let
$f_{i,\omega}\in(E^a)^*$ be a collection of functionals annihilating its
image.  The loop columns first imply that $f_{i,\omega}$ annihilates
$L_i$.  Write its induced functional on $H_i$ uniquely as
\[
 f_{i,\omega}=p_i(c_{i,\omega}),\qquad c_{i,\omega}\in G.
\]
Here and below we use the same notation for the induced functional and
its pullback to $E^a$.  By
\eqref{eq:correction-paired-annihilator} and
\eqref{eq:correction-leading-pair}, annihilation of the reverse columns
is equivalent to
\begin{equation}
 c_{i,s_cB}=c_{j,Bs_c}
 \quad\text{on every edge }i\leftarrow j\text{ of color }c.
 \label{eq:correction-edge-equality}
\end{equation}
The two offsets of either color are consecutive integers, say $a,a+1$.
Applying \eqref{eq:correction-edge-equality} to both edges out of $j$
gives
\[
 c_{j+a,s_cB}=c_{j+a+1,s_cB}.
\]
As $j$ ranges over $\mathbb Z/11\mathbb Z$, this makes the value
independent of the vertex.  Equation~\eqref{eq:correction-edge-equality}
then identifies its value with that on the rotated word $Bs_c$.
Iterating rotations shows that the annihilator consists exactly of one
element of $G$ per cyclic orbit of words.  The converse follows from
the same paired-annihilator identity.

Equation~\eqref{eq:correction-cyclic-input} says that every such
annihilator kills the degree-$\ell$ error.  Finite-dimensional linear
duality therefore puts the negative error in the image of the correction
map.  Choose a correction solving all selected equations.  Lower degrees
remain solved, so transfer and the active partition imply that every
nonselected equation vanishes in degree $\ell$ as well.

Continue inductively.  A coefficient of degree $n$ can change only at
steps $n$ and $n+1$, so the resulting free series is well-defined.
Every fixed Schur coefficient is eventually zero and remains zero.
The construction also covers $g=0$, when the annihilator is zero.
For a periodic word the orbit consists of its distinct rotations; no
multiplicity is inserted into the annihilator calculation or into
\eqref{eq:correction-cyclic-input}.
\end{proof}

\begin{figure}[tbp]
\centering
\begin{tikzpicture}[
  x=1cm,y=1cm,>=Stealth,
  every node/.style={font=\small},
  vertex/.style={circle,draw=black!65,fill=white,minimum size=5.8mm,inner sep=1.5pt},
  block/.style={rectangle,rounded corners=2pt,draw=black!65,fill=white,
    minimum width=12mm,minimum height=6.5mm,inner sep=3pt},
  first/.style={draw=blue!65!black,line width=.85pt,->},
  second/.style={draw=orange!80!black,line width=.85pt,->},
  correction/.style={draw=black!60,line width=.75pt,dashed,->},
  arrowlabel/.style={fill=white,inner sep=2pt}
]
  \node[font=\small\bfseries] at (2.45,5.55) {The offset-one cycle};
  \foreach \j in {0,...,10}{
    \node[vertex] (v\j) at ({2.45+2.1*cos(90-360*\j/11)},
                             {2.9+2.1*sin(90-360*\j/11)}) {$\j$};
  }
  \foreach \j in {0,...,10}{
    \pgfmathtruncatemacro{\next}{mod(\j+1,11)}
    \draw[first] (v\j) -- (v\next);
  }
  \node[align=center,text width=27mm] at (2.45,2.9)
    {$j\longrightarrow j+1$\\[3pt]first letter $s_1$};
  \node[font=\small\bfseries] at (9.65,5.55) {All four offsets at one vertex};
  \node[block] (j) at (7.25,2.9) {$j$};
  \node[block] (j1) at (11.85,4.65) {$j+1$};
  \node[block] (j2) at (11.85,3.48) {$j+2$};
  \node[block] (j3) at (11.85,2.32) {$j+3$};
  \node[block] (j4) at (11.85,1.15) {$j+4$};
  \draw[first] (j) -- node[arrowlabel,pos=.64,above,sloped] {$s_1$} (j1);
  \draw[first] (j) -- node[arrowlabel,pos=.64,above,sloped] {$s_1$} (j2);
  \draw[second] (j) -- node[arrowlabel,pos=.64,below,sloped] {$s_2$} (j3);
  \draw[second] (j) -- node[arrowlabel,pos=.64,below,sloped] {$s_2$} (j4);
  \node[align=center] at (9.55,.45) {Translate this pattern modulo $11$.};
  \node[block] (pairj) at (3.2,-1.1) {$j$};
  \node[block] (pairi) at (9.8,-1.1) {$i$};
  \draw[first] (pairj.north east) to[bend left=12]
    node[arrowlabel,above] {$U s_c$: forward term} (pairi.north west);
  \draw[correction] (pairi.south west) to[bend left=12]
    node[arrowlabel,below] {$V B$: reverse correction, $|B|=\ell-1$} (pairj.south east);
\end{tikzpicture}

\caption{The eleven-block input pattern. Indices lie in $\mathbb Z/11\mathbb Z$.
The left panel shows the offset-one cycle; the right panel specifies all
four forward offsets, translated at every vertex. Each forward arrow carries a
generic pair of input matrices. At order $\ell\ge3$, a reverse correction
has degree $\ell-1$; degree-$\ell$ loop corrections are not drawn.}
\label{fig:eleven-blocks}
\end{figure}

Figure~\ref{fig:eleven-blocks} separates the generating cycle from the
full local arrow pattern. The role of the four offsets is now explicit.  Avoiding returns of length
one or two starts the induction without changing the linear data.
Consecutive offsets within each color make its only obstructions
independent of the vertex.  The remaining obstructions are exactly the
cyclic sums already proved to vanish.

\subsection{Fixed roots and Schur complements}

The following elementary fact identifies the spectral equations with a
single polynomial identity. Its coefficient ring need not be commutative.
Let $E$ be a field, let $\mathfrak a=(s_1,s_2)$ be the augmentation ideal
of the free algebra $E\langle s_1,s_2\rangle$, and put
\[
  T_N=E\langle s_1,s_2\rangle/\mathfrak a^{N+1}
  \qquad (N\geq 0).
\]
Elements of $E$ are central in $T_N$. Reduction modulo $\mathfrak a$
is called the constant term.

\begin{lemma}\label{lem:schur-polynomial}
Let $W\in M_d(T_N)$ have constant term
\[
  W_0=\sum_{h=1}^{e}\lambda_h\pi_h,
  \qquad
  \pi_h\pi_j=\delta_{hj}\pi_h,
  \qquad
  \sum_{h=1}^{e}\pi_h=I_d,
\]
where the $\lambda_h\in E$ are distinct and the $\pi_h\in M_d(E)$
are nonzero projections. Set
\[
  q(t)=\prod_{h=1}^{e}(t-\lambda_h),
  \qquad \bar\pi_h=I_d-\pi_h.
\]
On the complementary summand $\bar\pi_h T_N^d$, the operator
\[
  C_h=\bar\pi_h(W-\lambda_h I_d)\bar\pi_h
\]
is invertible. Define the Schur expression on $\pi_h T_N^d$ by
\[
  S_h=\pi_h(W-\lambda_h I_d)\pi_h
      -\pi_hW\bar\pi_h C_h^{-1}\bar\pi_hW\pi_h.
\]
Then $q(W)=0$ if and only if $S_h=0$ for every $h$.
\end{lemma}

\begin{proof}
The constant term of $C_h$ is
$\sum_{j\ne h}(\lambda_j-\lambda_h)\pi_j$, which is invertible on
$\bar\pi_h E^d$. A matrix over $T_N$ with invertible constant term is
invertible: after factoring out its constant term, the inverse is a
finite geometric series in a matrix with entries in $\mathfrak a$.
This proves the assertion about $C_h$. All modules below are right
$T_N$-modules, with matrices acting on column vectors from the left.

Suppose first that every $S_h$ vanishes. Define
$K_h:\pi_hT_N^d\longrightarrow T_N^d$ by
\[
  K_hu=u-C_h^{-1}\bar\pi_hW\pi_hu,
\]
where the second term lies in the complementary summand. Then
$\pi_hK_hu=u$, and direct block multiplication gives
\[
  (W-\lambda_h I_d)K_hu=S_hu=0.
\]
The map $K=\sum_hK_h\pi_h$ has constant term $I_d$, so it is
invertible. The displayed kernel equations imply
$WK=KW_0$. Consequently
\[
  q(W)=Kq(W_0)K^{-1}=0.
\]

Conversely, suppose $q(W)=0$. Define
\[
  P_h=\prod_{j\ne h}
       \frac{W-\lambda_j I_d}{\lambda_h-\lambda_j}.
\]
These products are unambiguous because their factors are polynomials
in the same matrix with central scalar coefficients. Polynomial
interpolation modulo $q$ shows that the $P_h$ are spectral
idempotents; in particular,
\[
  (W-\lambda_h I_d)P_h=0,
  \qquad P_h\equiv\pi_h\pmod{\mathfrak a}.
\]
Let $\iota_h:\pi_hT_N^d\longrightarrow T_N^d$ be inclusion. The
endomorphism $A_h=\pi_hP_h\iota_h$ of $\pi_hT_N^d$ has constant
term the identity, and is therefore invertible. Thus
\[
  \widetilde K_h=P_h\iota_hA_h^{-1}
\]
satisfies
$\pi_h\widetilde K_h=\operatorname{id}_{\pi_hT_N^d}$ and
$(W-\lambda_h I_d)\widetilde K_h=0$. The complementary component
of the latter equation determines its remaining component uniquely:
\[
  \bar\pi_h\widetilde K_h
    =-C_h^{-1}\bar\pi_hW\pi_h.
\]
Its $\pi_h$ component now gives $S_h=0$.
\end{proof}

For a matrix over $E\langle\!\langle s_1,s_2\rangle\!\rangle$,
apply Lemma~\ref{lem:schur-polynomial} after reduction to $T_N$.
It follows that the fixed-root Schur equations hold through degree
$N$ exactly when the identity $q(W)=0$ holds through degree $N$.
In particular, this equivalence applies at every finite stage of a
formal correction, with the same polynomial $q$ throughout.
